\scanpage{25}
\mspage{12}

(On finit par avoir $I_{\alpha_0}=I$, bien sûr.) Les \deleted{\unsure{pose}}

\[
\begin{aligned}
 X_\alpha&=X_{I_\alpha}=\underset{i\in I_\alpha}{\underrightarrow{\lim}}\,X_i
       &&\unsure{u}\\
 Y_\alpha&=Y_{I_\alpha}=\underset{i\in I_\alpha}{\underrightarrow{\lim}}\,Y_i\\
 u_\alpha&=u_{I_\alpha}:X_\alpha\mathbin{\unsure{-|\to}}Y_\alpha,
\end{aligned}
\]

\unsure{forment|hom} \unsure{des|de} systèmes inductifs transfinis. On aura

\[
 \deleted{\unsure{u}}X_{\alpha_0}=X_I,\ Y_{\alpha_0}=Y_I,\ u_{\alpha_0}=u_I,
\]

et on doit \inserted{donc} prouver que $u_{\alpha_0}\in C$. Pour ceci on applique la prop.~1, \deleted{\unsure{elle}} \deleted{\unsure{il nous|on aura}} \inserted{suffit} donc de prouver que $(u_\alpha)\in C_{I_{\alpha_0+1}}$,

\[
 [\,\text{où }I_{\alpha_0+1}=\{\text{ordinaux }i\mid i<\alpha_0+1\ \text{i.e. }i\leq\alpha_0\}\,]
\]

\deleted{est dans $C$} Donc il faut prouver que les flèches $\widetilde u_{i_0}$ du cas p.~1 (où cette fois les $i,i_0$ sont des ordinaux, $i_0\leq\alpha_0$) sont dans $C$. Je vais distinguer les deux cas où $\alpha=i_0$ est ordinal \deleted{\illegible{1}} limite, \deleted{\unsure{qui est|quelque} \unsure{trivial}} \inserted{trivial car alors} $X_{i_0}\xleftarrow{\sim}\underset{i<i_0}{\lim}X_i$, $Y_{i_0}\xleftarrow{\sim}\underset{i<i_0}{\underrightarrow{\lim}}Y_i$, et $\widetilde u_{i_0}=\mathrm{id}$, et le cas $\alpha$ non \deleted{\unsure{limite}}, \unsure{i.e.|ici}.

\diagram{\begin{tikzcd}[ampersand replacement=\&,column sep=huge]
 X_{I_{\alpha-1}}\arrow[r]\arrow[d,"u_{\alpha-1}"']\&X_{I_\alpha}\arrow[d,"\beta_\alpha"]\arrow[dd,bend left=46,"u_\alpha"]\\
 Y_{I_{\alpha-1}}\arrow[r,no head]\&\widetilde Y_{I_\alpha}\arrow[d,"\widetilde u_\alpha"]\\
 \&Y_{I_\alpha}
\end{tikzcd}}
% archive sheet number: 24

